\documentclass[a4paper,12pt,oneside]{book} % A4纸，小四号大小就是12pt
\usepackage[utf8]{inputenc}
\usepackage{graphicx}

% 字体相关设置
\usepackage[AutoFakeBold=true]{xeCJK} %使用CJK包以便使用中文
\usepackage{CJKnumb}
\usepackage{xunicode}
\usepackage{fontspec}
\setmainfont{Times New Roman}
\setCJKmainfont{SimSun}
\setCJKfamilyfont{Heiti}{SimHei}

% 字体大小
\newcommand{\wuhao}{\fontsize{10pt}{10pt}\selectfont} % 五号字体 10pt
\newcommand{\xiaoer}{\fontsize{18pt}{0pt}\selectfont} % 小二字体 18pt
\newcommand{\sanhao}{\fontsize{16pt}{0pt}\selectfont} % 三号字体 16pt

% 页面设置
\usepackage[headheight=2cm,marginparwidth=0pt,headsep=1cm,foot=1cm,marginparsep=0pt]{geometry} % 设置页眉高度2cm,页脚高度其实没卵用的，除了上边距+页眉高度+正文体高度，其他都是页脚
\geometry{left=2.5cm,right=2.5cm,top=3cm,bottom=2.5cm} % 设置左、右、下边距为25mm，上边距30mm

% 页眉和页脚设置
\usepackage{fancyhdr} % 需要使用的包
\fancypagestyle{plain}{% 重定义 plain style，这个样式章节中全部用到。
	\fancyhf{}
	\fancyhead[C]{\wuhao 湖南科技大学本科毕业设计(论文)}
	\fancyfoot[C]{-\thepage-}
%	\renewcommand{\headrulewidth}{0pt} % 去掉页眉的线
}

% 章节设置
\usepackage{titlesec}
	% 设置 章 标题
\titlespacing*{\chapter}{0pt}{18pt}{18pt} % 段前段后各空0.5行
\titleformat{\chapter}[hang]{\centering\xiaoer\bfseries}{第\CJKnumber\thechapter 章\quad}{0pt}{}[\vspace{0cm}]
	% 设置 节 标题
\newcommand{\sihao}{\fontsize{14pt}{0pt}\selectfont} % 四号字体14pt
\titlespacing*{\section}{0pt}{14pt}{14pt}% 段前段后各空0.5行
\titleformat{\section}[hang]{\sihao\bfseries}{\thesection\quad}{0pt}{}{}
	% 设置 条 标题
\newcommand{\xiaosi}{\fontsize{12pt}{0pt}\selectfont} % 小四字体12pt，空0.5行
\titlespacing*{\subsection}{0pt}{12pt}{12pt}
\titleformat{\subsection}[hang]{\xiaosi\bfseries}{\thesubsection\quad}{0pt}{}{}

% 使用款、项
\usepackage{enumitem}
	% label=(\arabic*) 款、项序号为 (1), (2)...
	% itemindent=24 一个字12pt，空两格24pt
	% itemsep=-3pt 每两个item的上下间距减小，更紧凑
\setenumerate{label=(\arabic*), itemindent=24pt, itemsep=-3pt}

% 正文设置
\usepackage{indentfirst} % 首行缩进
\setlength\parindent{24pt} % 因为字大小12pt，所以缩进两个字就是24pt
\linespread{1.3} % 1.25倍行距 [1是单倍行距，1.3是1.25倍行距，1.6则是双倍行距]
\usepackage{layout}

% 目录设置
\usepackage{titletoc} % 控制目录的宏包
\titlecontents{chapter}[0em]{\addvspace{18pt}\sihao}{\bfseries 第\CJKnumber\thecontentslabel 章\quad}{}{\wuhao\titlerule*[5pt]{$\cdot$}\contentspage} % 目录章样式，四号加粗宋体 titlerule*[Xpt]调整点的密集度，X越小密度越高
\titlecontents{section}[2em]{\addvspace{14pt}\xiaosi}{\thecontentslabel\quad}{}{\wuhao\titlerule*[5pt]{$\cdot$}\contentspage}% 目录节样式
\titlecontents{subsection}[4em]{\addvspace{8pt}\xiaosi}{\thecontentslabel\quad}{}{\wuhao\titlerule*[5pt]{$\cdot$}\contentspage}
\renewcommand{\contentsname}{\sanhao\CJKfamily{Heiti}\textbf{目 \qquad 录}} % 目录样式

% 脚注设置
\usepackage{tikz}
\newcommand{\footnotenumber}{\fontsize{4pt}{0pt}\selectfont}
\newcommand*\circled[1]{%
	\tikz[baseline=-4pt]\node[shape=circle,draw,inner sep=0.2pt,font=\footnotenumber,minimum size=1pt] (char) {#1};} % baseline=X，X负的绝对值越大，脚注产生的圆圈就越高
\renewcommand\thefootnote{\protect\circled{\arabic{footnote}}}

% 数学相关设置
\usepackage{amsbsy}
\usepackage{amssymb}
\usepackage{amstext}

% 插图编排设置
\usepackage{float} % 能够使用 H 位置
\usepackage{wrapfig} % 能够使用文字环包的图
\usepackage{graphicx}
\usepackage[labelfont=bf,font=bf]{caption}
\usepackage{subcaption} % 子图的标题样式需要这个包
\captionsetup[subfigure]{labelformat=simple, labelsep=none}
\renewcommand{\thesubfigure}{\alph{subfigure})\ } % 子图的标题样式 a) b) c)

% 参考文献
\renewcommand{\bibname}{\sihao 参考文献}
\bibliographystyle{gbt7714-2005} % 使用 GB/T 7714-2005 的标准，GB7714-87在2005年就被废了 感谢提供：https://github.com/Haixing-Hu/GBT7714-2005-BibTeX-Style
\usepackage{url} % 使用在线链接参考文献使用

% 致谢设置
\newcommand{\thanksfor}{\newpage\addcontentsline{toc}{chapter}{\sihao \bfseries 致\qquad 谢}\chapter*{致\qquad 谢}}

% 附录设置
\newcounter{appnum}
\setcounter{appnum}{0}
\renewcommand{\theappnum}{\Alph{appnum}}
\newcommand{\NewAppend}{\refstepcounter{appnum}\newpage\addcontentsline{toc}{chapter}{\sihao \bfseries 附\quad 录 \Alph{appnum}}\chapter*{附\quad 录 \Alph{appnum}}}

% 图表label设置
\usepackage{multirow} % 多行列表
\renewcommand{\figurename}{图}
\renewcommand{\tablename}{表}

\usepackage{pdfpages} % 将扉页、任务书等PDF进行导入需要用到此包

\usepackage{listings} % 插入代码需要

\begin{document}

\author{author}
\title{title}
\date{2016年上半年}

\frontmatter
\includepdf[pages={1-}]{../docs/frontpages.pdf} % 将扉页、任务书等PDF导入，如果不需要，可以注释这行。
%\maketitle

% 设置摘要，因为 class 是 book，所以没有 Abstruct 环境，而且字体大小跟chapter一样，所以用 chapter* 来代替
\newcommand{\NewAbstract}[1]{\chapter*{#1}\thispagestyle{empty}}
\newpage
\pagestyle{empty}
\NewAbstract{\xiaoer\bfseries 摘\qquad 要}

中文摘要写在这。

\bigskip % 空行
{\CJKfamily{Heiti}\bfseries 关键词：}
 % 请转到该文档中填写中文摘要
\NewAbstract{\xiaoer\bfseries ABSTRACT}

English Abstract here

\bigskip % Blank Line
{\CJKfamily{Heiti}\bfseries KeyWords:\quad} % 请转到该文档中填写英文摘要

\newpage
\setcounter{page}{1}
\pagestyle{plain} % 目录开始使用统一的页面样式

% 插入目录
\tableofcontents

\mainmatter
\newpage

% 以下是各个章：
% 每章可自定义名称，然后包含进来
\chapter{示例章节}
示例文字：跟着我左手右手一个慢动作
\section{示例小节}
示例文字：右手左右慢动作重复播
\subsection{示例条}
示例文字：这首歌\\ 给你快乐\\ 你有没有爱上我\footnote{tfboys的歌，你还好吗？}，哈哈

\begin{enumerate}[leftmargin=24pt]
	\item 昨天所有的荣耀
	\item 已变成遥远的回忆
	\item 辛辛苦苦，已度过半生
	\item 今夜重又走进风雨
\end{enumerate}

\begin{figure}[H]
	\centering
	\begin{subfigure}[b]{0.3\textwidth}
		\centering
		\includegraphics[width=.15\textheight]{images/make_a_big_news.jpg}
		\caption{弄个大新闻}
	\end{subfigure}
	\begin{subfigure}[b]{0.3\textwidth}
		\centering
		\includegraphics[width=.15\textheight]{images/you_like_big_news.jpg}
		\caption{不要想喜欢弄大新闻}
	\end{subfigure}
	\caption{我为长者续一秒}
\end{figure}

{\wuhao 五号字体}

{\CJKfamily{Heiti} 然后搞点黑体}

举个文献的栗子\footnote{感谢提供：https://github.com/Haixing-Hu/GBT7714-2005-BibTeX-Style}：诸葛亮写了本书叫《论六出祁山的必要性与可行性》\textsuperscript{\cite{诸葛亮:Book}}，然后发了篇文章\cite{诸葛亮:Article}，然后刘备就长坂坡摔小孩了\textsuperscript{\cite{刘备:摔小孩}}\footnote{例子原型：http://lyanry.is-programmer.com/posts/195.html}

\subsection{又来一节}
你的头发还好好吗？


\backmatter
% bibliography, glossary and index would go here.

% 参考文献
\newpage
\addcontentsline{toc}{chapter}{\sihao \bfseries 参考文献} % 把参考文献加入目录
{\wuhao \bibliography{bibfile}} % 请在 bibfile.bbl 中写好文献 % 要求用5号字体

% 致谢
% 致谢请用 \thanksfor
\thanksfor

% 附录
% 附录请用 \NewAppend
\appendix
\NewAppend
这是个附录

\end{document}
